\documentclass{article}
\usepackage{tackling_climate_workshop_style}
\usepackage[utf8]{inputenc}
\usepackage[T1]{fontenc}
\usepackage{microtype}
\usepackage{booktabs}
\usepackage{amsmath}
\usepackage{graphicx}
\usepackage{xcolor}
\usepackage{hyperref}
\usepackage{url}
\hypersetup{hidelinks}

\newcommand{\system}{ClimateKG}
\newcommand{\code}[1]{\texttt{#1}}
\title{ClimateKG: Context-Gated Scientific GraphRAG for\\Literature-Grounded Land-Use Scenario Screening}
\author{Anonymous Author(s)}

\begin{document}
\maketitle

\begin{abstract}
Land-use and land-cover changes alter surface energy, water, and momentum exchanges, but the resulting cloud, rainfall, and temperature responses depend on season, hydroclimate, background wind, and intervention geometry. Testing many spatial land-use scenarios with convection-resolving models is expensive, while topical retrieval can combine environmentally incompatible findings. We introduce \system{}, a scientific GraphRAG system that extracts evidence-backed claims with the complete study settings under which they hold, gates claims by a query's environmental context before graph traversal, and returns answers traceable to exact source passages. Our shared pilot graph contains 18 finalized papers, 181 study contexts, and 279 scientific claims. A ten-paper structural audit found that every extracted scope and evidence reference resolved; a ten-question comparison with passage RAG and community-based graph retrieval measures claim recall, condition retention, and citation resolution. Case studies show why conditioning matters: patch size and wind alter heterogeneous-surface convection, while one deforestation pattern yields null, cloud-only, or rainfall responses across atmospheric regimes. \system{} is intended to propose auditable hypotheses and candidate land-use rasters for subsequent WRF evaluation, not to prescribe land management or replace simulation.
\end{abstract}

\section{Introduction}
Land surfaces influence temperature, boundary-layer growth, clouds, rainfall, and circulation through albedo, roughness, soil moisture, and evapotranspiration \citep{pielke2007,cao2020}. These effects are not intervention lookup rules. Mesoscale soil-moisture patterns can organize storm initiation \citep{taylor2011}; changing patch size or background wind can change whether and where convection develops \citep{lee2019}. Likewise, observed Amazonian deforestation produced different cloud and rainfall responses under different seasonal and synoptic regimes \citep{wang2000}. A useful literature system must therefore retrieve not only \emph{what} changed and \emph{what} followed, but \emph{where, when, at what scale, and under which atmospheric state}.

The motivating application is scenario design. Exploring many high-resolution land-cover layouts for many regions with a numerical model such as WRF \citep{skamarock2019} is computationally demanding; multi-year convection-resolving regional simulations illustrate the scale of this cost \citep{leutwyler2017}. Literature can narrow this search: retrieve transferable mechanisms, identify conditions and counter-evidence, generate a small set of spatially explicit candidate land-use rasters, and then test those candidates numerically. Plain RAG \citep{lewis2020} provides passage evidence, and generic GraphRAG supports corpus-level connections and summaries \citep{edge2024}; neither representation makes environmental applicability a first-class constraint.

We contribute (1) a provenance-preserving representation of scientific findings and their study settings; (2) context gating before mechanism-path search, while separating missing information from mismatch; (3) staged and single-call extraction routes for different paper lengths; and (4) a reproducible pilot with saved prompts, model responses, intermediate JSON, graph tables, query reports, and exact citations. The contribution is an evidence-to-hypothesis system, not evidence that any retrieved intervention will work locally.

\begin{figure}[t]
\centering
\fbox{\begin{minipage}{0.96\linewidth}\small
\textbf{Indexing:} PDF $\rightarrow$ Nemotron Parse v1.2 $\rightarrow$ cleaned \textbf{EvidenceBlocks} $\rightarrow$ hierarchical \textbf{StudyContexts} + \textbf{ContextFacets} $\rightarrow$ \textbf{LandTransitions} + evidence-backed \textbf{ScientificClaims} $\rightarrow$ canonical concepts + embeddings $\rightarrow$ Parquet/Neo4j graph.\\[-1pt]
\textbf{Querying:} question + optional local data $\rightarrow$ structured query context $\rightarrow$ context candidates and coverage $\rightarrow$ claim gate $\rightarrow$ compatible mechanism paths $\rightarrow$ smallest supporting EvidenceBlocks $\rightarrow$ cited answer.
\end{minipage}}
\caption{\system{} makes the scientific setting part of retrieval. Claims are gated before they can participate in a mechanism path.}
\label{fig:pipeline}
\end{figure}

\section{Methodology}
\paragraph{Literature corpus and comparison protocol.}
The working collection contains 26 land--atmosphere PDFs; 18 finalized, deduplicated papers form the shared corpus. All systems use the same cleaned text, Qwen3.6--27B generator, Qwen3-Embedding-4B vectors, and top-eight evidence budget. Baseline synthesis uses medium reasoning; \system{} uses its validated stage profile (no reasoning for parsing and low/medium for synthesis). PlainRAG retrieves passages. The graph baseline applies Microsoft GraphRAG 3.1.1 hierarchical Leiden clustering to canonical State entities and ScientificClaim relationships, then reranks relationships from five communities. It tests community retrieval over the same extraction, not Microsoft GraphRAG's raw-text extractor. \system{} additionally gates Claims by StudyContext before path search.
\paragraph{Scientific representation.}
An \textbf{EvidenceBlock} is a stable page-linked passage. A hierarchical \textbf{StudyContext} represents one complete setting while sharing common setup across seasonal, spatial, or simulation children. A \textbf{ContextFacet} records one environmental property; a \textbf{LandTransition} links baseline and changed settings. A \textbf{ScientificClaim} connects canonical source and target States, distinguishes causal from associative language, and cites the smallest supporting blocks.

\paragraph{Indexing.}
NVIDIA Nemotron Parse v1.2 is the sole PDF parser. Deterministic cleaning removes repeated furniture and excludes references from semantic extraction while preserving tables, units, and provenance. Long papers use staged section scouting, a compact hierarchy of complete settings, facet/claim extraction, consolidation, and State canonicalization. Short cleaned papers may use one structured extraction call before the same deterministic downstream stages. Schemas reject unresolved references; embedding similarity proposes but does not force concept merges.

\paragraph{Context retrieval and claim gating.}
The query parser returns optional source/target States, an intervention, and ContextFacets. Approximate nearest-neighbor search proposes contexts; exact same-domain facet MaxSim then reranks them with a coverage penalty. Missing facets reduce coverage but are not mismatches. A Claim inherits applicability from its supporting context, while only materially conditioning facets affect its Claim score. Context, intervention, mechanism, endpoint, and spatial-configuration signals remain separately auditable. Gating precedes beam search, so a path cannot join Claims from incompatible settings. Relation-specific spatial translation additionally requires the corresponding local geometry.

\paragraph{Grounded answer.}
The synthesis prompt sees selected claim records, context scores, and SourceBlock references; the application retains the full passages for provenance and citation rendering. Evidence-budget trimming cannot retain a claim after removing all its support. The structured response must attach the smallest supporting Claim-ID set to each statement, which is then rendered as paper, page, Claim, and block citations. Scores are ranking signals, not probabilities that a finding is true.

\section{Results}
\paragraph{Indexing and structural audit.}
The 18-paper comparison corpus contains 2,746 EvidenceBlocks, 181 StudyContexts, 556 ContextFacets, 58 LandTransitions, and 279 ScientificClaims. Canonicalization produces 399 unique State nodes linked by 279 Claim edges. The graph has 132 weakly connected components; its largest contains 59 States, while seven States support 16 possible directed two-hop Claim pairs. The graph therefore contains mechanism chains but remains sparse, which motivates retaining direct evidence alongside paths.

Across ten papers selected by intervention and study type, we audited 1,539 blocks, 114 contexts, 323 facets, 35 transitions, and 185 claims. Every Claim scope, conditioning Facet, and evidence reference resolved. This establishes structural validity, not scientific recall: independent experts have not yet enumerated all findings, so we leave recall unreported rather than using the extraction as its own ground truth.

\paragraph{Ten-question retrieval audit.}
We fixed ten questions before the comparison: nine in-corpus questions test scale, wind, seasonal regimes, observational/model disagreement, edge effects, sign differences, remote effects, hydroclimate reversals, and counteracting mechanisms; one out-of-corpus Rajasthan question tests cautious transfer. Expected Claim IDs and short reference propositions were prepared from the indexed corpus. They are author annotations, not independent gold labels. Graph systems are scored by exact Claim ID; PlainRAG passages are mapped to same-paper Claims using the shared embedding model (cosine $\geq0.55$) or direct SourceBlock term coverage ($\geq0.75$). We separately measure retrieved-Claim recall, expected Claims used in the cited answer, whether citations resolve, sentence citation coverage, and retention of requested conditions. Full prompts, thinking fields, retrieval records, and answers are saved.

The first completed cases distinguish the systems' behavior. For the zero-wind 14.4-km question, all three recovered the four expected Claims and produced resolvable citations. PlainRAG nevertheless retrieved four of eight passages from other studies; its answer remained on the correct LES evidence. The community graph recovered all four expected Claims but added uncited real-world caveats not supplied in its evidence. \system{} used all four and explicitly blocked spatial translation. For the 1~m~s$^{-1}$ question, PlainRAG supported two of five expected Claims and correctly reported missing convergence/cloud-motion evidence; community retrieval recovered four, while \system{} used all five and preserved their zero-to-one-wind comparison. These are useful failure traces, not sufficient evidence of general superiority.

Table~\ref{tab:cases} presents two manually checked cases. For the in-corpus heterogeneous-surface query, parsing preserved the 14.4-km chessboard and zero-wind facets. The top context had full facet coverage ($S_{ctx}=0.676$), and the answer cited the finding that patches above 5 km transitioned from shallow cumulus to deep precipitating convection over dry patches, while homogeneous forcing retained shallow clouds. It also preserved the counterintuitive result that cross-patch circulation reduced the near-surface dry--wet temperature contrast \citep{lee2019}. With 1 m s$^{-1}$ wind, a separate claim moves convergence from the dry-patch center to its downwind edge, so collapsing the cases would produce a misleading spatial answer.

The Rondonia extraction separates rainy, break, and dry-season simulations. Under rainy-season strong synoptic forcing, control and deforestation rainfall were nearly identical and mesoscale circulations were inhibited. During the break, weakened winds allowed enhanced low cloud but little colocated rainfall. In the dry-season setting, weak wind allowed a stronger land-induced circulation and localized convective response \citep{wang2000}. The same land pattern therefore cannot be mapped to one unconditional ``deforestation increases rainfall'' edge.

\begin{table}[t]
\centering\scriptsize
\caption{Why complete study settings matter. IDs point to released claims and exact supporting blocks in the appendix.}
\label{tab:cases}
\begin{tabular}{p{0.18\linewidth}p{0.35\linewidth}p{0.35\linewidth}}
\toprule
Case & Preserved conditions & Indexed/queried conclusion \\
\midrule
Heterogeneous LES & 14.4-km dry/wet chessboard; zero wind & $>5$-km patches support deep precipitating convection over dry patches; cross-patch mixing reduces surface contrast (CL002, CL004). \\
Rondonia deforestation & Rainy/break/dry season; original vs. weakened synoptic wind & Null rainfall under strong forcing; cloud-only response in weakened-wind break case; dry-season circulation response (CL001, CL006--010, CL015). \\
New Rajasthan location & Semi-arid; 10--15-km alternating irrigated/dry patches; weak monsoon wind & Literature analogy can motivate a mechanism hypothesis, but cannot establish local rainfall or cardinal placement without local geometry and numerical/field validation. \\
\bottomrule
\end{tabular}
\end{table}

\paragraph{Extraction routes and cost.}
An archived staged run completed 17 of 20 registered papers and made 305 semantic calls, averaging 20.02 minutes per completed paper. A separate combined-extraction experiment made one semantic call per cleaned paper; 7 of 10 outputs passed all checks and successful papers averaged 3.29 minutes with cached parsing. The rejected combined outputs contained invalid evidence or conditioning references. This supports routing short cleaned papers to combined extraction, but also shows why the staged hierarchy and fail-closed checks remain necessary for complex papers. Route times are not directly comparable because parsing was cached only in the combined experiment.

\paragraph{Interpretation.}
The expected Claims are derived from the same extraction that supplies the graph baselines, so the retrieval audit tests whether each query system recovers and uses known indexed evidence; it does not evaluate whether indexing missed findings from the PDFs. Citation resolution likewise does not establish entailment. The full ten-query table is generated only after every saved trace completes, and independent paper-level annotation and clause-level citation review remain required for a defensible accuracy claim.

\section{Climate impact pathway and limitations}
For scientists, \system{} exposes mechanism paths and contradictions for literature synthesis. For regional planners, it can translate a local question into a shortlist of evidence-backed scenarios and explicit unknowns. For farmers or communities, it may support transparent discussion, but should not issue autonomous recommendations. A proposed workflow is: define the target region and objective; derive citable local climate, terrain, land-cover, wind, and geometry variables; retrieve compatible findings; generate a small, diverse set of LULC rasters; evaluate them with WRF or another appropriate model; and assess water, food, biodiversity, livelihood, and justice trade-offs with local experts and affected communities.

The present corpus is small and geographically biased; extraction still fails; similarity is not causal validity; and no WRF experiment, field validation, stakeholder study, or independent system-accuracy benchmark has been completed. The graph is sparse, and a useful answer may rely on direct Claims rather than a multi-hop mechanism. The new-location case is consequently a transferability audit, not a recommendation. Immediate priorities are independent context/claim annotation, clause-level citation audits, broader equal-budget baselines, and numerical testing of selected raster scenarios. The present evidence supports a narrower conclusion: representing complete study settings prevents distinct regimes from being flattened into one unconditional literature edge and makes generated hypotheses inspectable at their sources.

\clearpage
\bibliographystyle{plainnat}
\bibliography{references}

\clearpage
\appendix
\section{Reproducibility and extraction detail}

\subsection{Released artifacts}
The repository preserves the rendered/cleaned document, stable EvidenceBlocks, every completed LLM request and response, schema-validation failures, consolidated paper JSON, embeddings, Parquet graph tables and ANN indexes, and per-query reports. The two extraction routes converge on the same final schema. The staged route is:
\begin{quote}\small
section scout $\rightarrow$ hierarchy of complete experimental settings $\rightarrow$ ContextFacet extraction $\rightarrow$ ScientificClaim extraction $\rightarrow$ context reconciliation $\rightarrow$ paper consolidation $\rightarrow$ state canonicalization $\rightarrow$ embeddings.
\end{quote}
The combined route supplies all cleaned non-reference EvidenceBlocks to a single structured extraction prompt, then uses the same reconciliation, consolidation, state, embedding, and graph steps. The routing count uses the \code{o200k\_base} tokenizer on cleaned non-reference blocks; 8,000 tokens is a preferred target rather than a hard rejection boundary, with a configured 16,000-token safety limit for the experiment reported here.

\subsection{Central object schemas}
Below, \code{evidence\_block\_ids} must name directly supporting passages from the current paper. IDs are permanent after assignment.
\begin{verbatim}
StudyContext {
  id, paper_id, parent_ids[], label, aliases[], spatial_support?,
  evidence_block_ids[], retrieval_embedding
}
ContextFacet {
  id, paper_id, context_id,
  domain: spatial_configuration | land_surface | hydrology |
          atmosphere | climate | substrate_terrain | other,
  notion, description, origin, source?, evidence_block_ids[], embeddings
}
LandTransition {
  id, paper_id, from_context_id, to_context_id,
  label, aliases[], description, evidence_block_ids[]
}
ScientificClaim {
  id, paper_id,
  scope_type: context | transition, scope_id,
  from: {concept, state}, to: {concept, state},
  relation: causal | associative,
  description, evidence_role,
  conditioning_facet_ids[], evidence_block_ids[], claim_embedding
}
\end{verbatim}

\section{Staged trace: heterogeneous surface experiment}
The indexed paper is Lee et al. \citep{lee2019}, DOI \code{10.1175/JAS-D-18-0196.1}. Nemotron Parse produced 163 EvidenceBlocks. Staged extraction produced 26 StudyContexts, 59 ContextFacets, four LandTransitions, and 18 ScientificClaims. The setting hierarchy retains the common LES setup and separates heterogeneous patch size/wind combinations from homogeneous dry, wet, and domain-mean controls. For example, the 14.4-km case has distinct zero, 1, and 2 m s$^{-1}$ children rather than one context with wind described in prose.

\paragraph{Claim CL002.}
\code{P000001\_CL002} is a causal claim scoped to transition \code{P000001\_T001}; it is conditioned by facet \code{P000001\_F012}. Its description states that chessboard heat-flux heterogeneity with patches above 5 km produces a transition from shallow cumulus to deep precipitating convection over dry patches, whereas homogeneous domain-mean fluxes retain scattered shallow clouds. The sole evidence is \code{P000001:S06:P0033}, which reports deep clouds above 6 km and significant precipitation over dry patches for large zero-wind patches, versus shallow cumulus for patches below 5 km.

\paragraph{Claim CL004.}
\code{P000001\_CL004} uses \code{P000001:S11:P0112} to record that active cross-patch circulation horizontally mixes the boundary layer and reduces the near-surface dry--wet temperature contrast. This is easy to omit from a simple ``heterogeneity causes convection'' summary, yet it is important because the surface-temperature response and cloud response do not follow the same naive contrast argument.

\paragraph{Wind-dependent spatial claim.}
\code{P000001\_CL014} cites \code{P000001:S11:P0097} and \code{P000001:S12:P0122}: changing background wind from zero to 1 m s$^{-1}$ moves the convergence/updraft from the dry-patch center to the downwind edge; sufficiently stronger flow destroys the organized circulation. This claim belongs to a different transition and cannot be substituted into the zero-wind answer.

\subsection{Query trace}
The frozen question was: ``Under zero background wind over a heterogeneous chessboard surface with 14.4 km dry and wet patches, how does surface heat-flux heterogeneity affect local cloud development and precipitation?'' The parser created two user-origin facets: (i) spatial configuration, notion \emph{chessboard surface pattern}, description \emph{14.4 km patches}; and (ii) atmosphere, notion \emph{zero background wind}. The top retrieved context matched both domains (coverage 1.0), with spatial and atmospheric pair scores 0.662 and 0.691 and overall context score 0.676. After claim gating, synthesis cited CL002 for deep precipitating convection, CL001 for stronger updrafts, CL003 for mixed-layer thermodynamics, and CL004 for reduced surface contrast. The response explicitly limited transfer to the described patch configuration and did not invent a local cardinal placement.

\section{Staged trace: Rondonia deforestation}
The indexed study is Wang et al. \citep{wang2000}. Its 142 EvidenceBlocks yielded 20 StudyContexts, 44 ContextFacets, eight LandTransitions, and 33 ScientificClaims. The hierarchy separates break-period original wind, break-period weakened wind, rainy-season strong forcing, and dry-season weak-wind settings.

\begin{table}[h]
\centering\small
\caption{Selected extracted findings and their exact evidence anchors.}
\begin{tabular}{p{0.12\linewidth}p{0.54\linewidth}p{0.20\linewidth}}
\toprule
Claim & Extracted result & EvidenceBlock(s) \\
\midrule
CL001 & Break period, original wind: at most 0.1 mm at isolated control locations and almost no rainfall in the deforestation run. & S08:P0056 \\
CL006--007 & Weakened wind: enhanced low cumulus at 900--800 hPa over deforested areas, but sparse rain cells did not overlap the enhanced cloud. & S08:P0062, P0064 \\
CL008--010 & Rainy season: control and deforestation rainfall/cloud distributions were nearly identical because stronger synoptic forcing inhibited mesoscale circulation. & S09:P0072, P0074; S11:P0101 \\
CL015 & Dry season, weak wind: significant low-level mesoscale heat flux indicated a land-induced circulation; related evidence reports localized convection/rainfall on the first day. & S10:P0091, P0092 \\
\bottomrule
\end{tabular}
\end{table}

This trace illustrates the representation's purpose. A graph with one unconditional \emph{deforestation $\rightarrow$ rainfall} edge would erase null and cloud-only regimes and could reverse the interpretation of a local query.

\subsection{Query redesign after an observed retrieval failure}
An earlier Rondonia run parsed the seasonal and wind facets and retrieved P000019 Claims as candidates, but State-seeded path search retained only generic deforestation paths. The resulting answer was grounded in the selected package yet omitted available regime-specific evidence. Root-cause review led to three minimal query changes: retain a direct-evidence lane after context gating; seed optional paths with semantically relevant Claims rather than only query endpoint States; and preserve spatial quantities such as 10--15 km as independently matchable Query Facets. The direct lane does not bypass context gating and does not force unrelated Claims into a mechanism chain. Regression tests and the current benchmark use this revised query flow.

\section{New-location transfer and baselines}
The frozen out-of-corpus question describes a semi-arid Rajasthan planning area, strong daytime heating, weak monsoon background wind, and proposed alternating irrigated/dry patches of 10--15 km. A place name alone is not used to infer climate. Query enrichment may add reproducibly sourced local facets (coordinates/geometry, climate class, aridity, elevation, wind, land cover, and patch metrics), with each derived value retaining dataset and operation provenance. Retrieval may then identify analogous patch-circulation mechanisms, but the answer must report missing or mismatched facets and cannot assert local rainfall, optimal LULC, or cardinal placement without applicable evidence and local numerical or field validation.

\paragraph{Observed out-of-corpus result.}
The query resolved Rajasthan to an approximate point and derived a TerraClimate aridity ratio of 0.209 (semi-arid), BWh K\"oppen--Geiger class, cropland land cover, terrain, and seasonal ERA5-Land winds with dataset provenance. End-to-end latency was 290.4 s: query parsing used 12.9 s (882 prompt and 325 generated tokens), while synthesis used 232.3 s (19,224 prompt and 1,026 generated tokens). The answer cited mechanisms for enhanced afternoon cloud over vegetated/cropland surfaces and downwind edge effects, and explicitly warned that the evidence came from other climate regimes and did not establish transfer to Rajasthan. However, the parser encoded the 10--15-km alternating pattern in the source state rather than a spatial-configuration facet, and the answer omitted that scale. It also failed to reconcile the user's weak-monsoon description with derived JJA wind of 3.09 m s$^{-1}$. The run therefore passed basic grounding/abstention checks but failed the full conditioning audit.

The release includes three runnable systems: the proposed pipeline, a project-owned PlainRAG baseline, and a community-retrieval adapter that uses Microsoft GraphRAG 3.1.1 hierarchical Leiden clustering over ClimateKG's State/Claim projection. A frozen ten-query file specifies the shared corpus, local generator, embedding model, evidence budget, expected Claims, and conditions. Complete human-readable traces contain every retrieved record, exact prompt, Qwen thinking field, and final response.

\section{Known limitations and next validation steps}
The pilot is not a representative literature census. PDF errors, extraction omissions, schema failures, canonicalization mistakes, and incomplete local enrichment remain possible. A Claim's retrieval score is not scientific confidence, external validity, feasibility, or net climate benefit. Before planning use, the system requires independently enumerated paper settings/findings, clause-level citation review, geographic coverage analysis, sensitivity analysis for context weights and thresholds, and WRF experiments for selected scenarios. Any intervention also requires water, biodiversity, food, livelihood, governance, and equity assessment with affected communities.

\end{document}
